\documentclass[conference]{IEEEtran}
\IEEEoverridecommandlockouts

% Revised from the author's draft "Cross-Language Static Analysis for ROS 2: Detecting Shadow
% Configuration Drift". The previous version is kept in rosfix_before_humanize.tex.
% Sections V-A to V-D report the author's measurements of an earlier rosfix version; all
% numbers there are unchanged. Sections V-E onwards were measured on the current code.

\usepackage{cite}
\usepackage{amsmath,amssymb}
\usepackage{graphicx}
\usepackage{booktabs}
\usepackage{tabularx}
\usepackage{xcolor}
\usepackage{listings}
\usepackage{algorithm}
\usepackage{algpseudocode}
\usepackage{url}
\usepackage[hidelinks]{hyperref}

\lstset{basicstyle=\ttfamily\footnotesize,columns=fullflexible,keepspaces=true,
        frame=single,framerule=0.3pt,aboveskip=4pt,belowskip=4pt}
\newcommand{\code}[1]{\texttt{#1}}

\begin{document}

\title{Cross-Language Static Analysis for ROS 2:\\Detecting Shadow Configuration Drift}

\author{\IEEEauthorblockN{Soham Kishor Walam}
\IEEEauthorblockA{Robotics and Autonomous Systems Laboratory\\
Advanced Engineering Research\\
Email: soham.walam2005@gmail.com}}

\maketitle

\begin{abstract}
A ROS 2 system is configured in four languages at once: XML launch files, YAML parameter files,
Python launch scripts and C++ node code. Existing analysis tools check each file type on its own,
so they miss errors that only appear when the files are read together. We call these errors
\emph{Shadow Configuration Drift} and present rosfix, a static analyzer that parses all four
languages with Tree-sitter and related parsers and checks them against one shared model. On
38,917 files from 12 ROS 2 repositories, rosfix reported 2,405 configuration issues with 87.3\%
precision on a manually validated sample. Of these issues, 13.0\% cannot be found by analyzing
any single file. The current version of rosfix resolves launch graphs the way the ROS 2 launch
system and \code{rcl} do, uses per-distribution knowledge packs for installed packages, and keeps
a repair only if re-analysis shows that it removes the finding. In a mutation benchmark on the
same 12 repositories, it detected all 161 injected faults at the mutated line. Its repair
assistant, a 289-weight ranker compiled into the binary, picks the intended value in 69.5\% of
held-out cases from unseen packages and names, against 54.1\% for lexical ranking, with 7
wrong edits in 1,791 cases against 30; across five package folds, however, it does not make
fewer wrong edits in every fold.
\end{abstract}

\begin{IEEEkeywords}
Static analysis, robotics middleware, configuration management, cross-language analysis
\end{IEEEkeywords}

% ===========================================================================
\section{Introduction}

ROS 2 is built on the Data Distribution Service (DDS) and ships with a security framework
(SROS2), and it is now used for safety-critical robots. Its configuration, however, is spread
over several languages. Launch files in XML or Python start nodes and pass them YAML parameter
files, and the nodes themselves are written in C++ or Python. Each layer can change what the
others set.

A common case: an architect sets \code{RELIABLE} QoS in a YAML file so that safety commands are
always delivered, while a developer hardcodes \code{.best\_effort()} in the C++ node to cut
latency. The node runs with best-effort QoS, and no tool that reads only the YAML file or only
the C++ file can see the conflict. We call such conflicts between intended configuration and
runtime behavior \emph{Shadow Configuration Drift}.

Existing ROS 2 tools each cover one file type. QoS Guard~\cite{lee2025qosguard} checks DDS XML
profiles, LiSA4ROS2~\cite{zanatta2024lisa4ros2} analyzes C++ source to generate security
policies, and HAROS~\cite{santos2021haros} computes code-quality metrics for ROS 1. None of
them checks that configuration stays consistent across files. The need is documented: Canelas
et al.\ found that 27 of 50 sub-categories of ROS misconfiguration were not addressed by any of
18 surveyed techniques~\cite{canelas2024misconfig}, and Chen et al.\ studied 121 interaction
bugs across ten ROS projects~\cite{chen2025ibugs}.

This paper makes the following contributions:
\begin{enumerate}
\item A formal definition of Shadow Configuration Drift (Section~\ref{sec:problem}).
\item rosfix, a cross-language static analyzer for ROS 2 configuration (Section~\ref{sec:design}).
\item An evaluation on 38,917 files from 12 repositories. It found 2,405 issues at 87.3\%
  precision, and 13.0\% of them (312) required cross-file analysis (Sections~\ref{sec:accuracy}
  and~\ref{sec:drift}).
\item A resolved launch graph that matches parameter-file keys the way \code{rcl} does and links
  each launched node to its source code (Section~\ref{sec:graph}).
\item Support for installed packages through the ament resource index and per-distribution
  knowledge packs (Section~\ref{sec:installed}).
\item A repair engine that re-analyzes the workspace before it keeps a fix, a repair assistant
  that only chooses among values valid at the site, and an optional sandboxed build check
  (Section~\ref{sec:repair}).
\item A mutation benchmark that injects faults at real sites in the 12 repositories. All 161
  faults were detected (Section~\ref{sec:mutation}).
\end{enumerate}

% ===========================================================================
\section{Problem Formulation}
\label{sec:problem}

\subsection{Shadow Configuration Drift}

Let $C = \{C_{xml}, C_{yaml}, C_{py}, C_{cpp}\}$ be the configuration sources of a ROS 2
system. The effective runtime configuration $C_{eff}$ is their composition under the precedence
rules:
\begin{equation}
C_{eff} = C_{cpp} \oplus C_{py} \oplus C_{yaml} \oplus C_{xml}
\end{equation}
where $\oplus$ is override composition (later sources override earlier ones).

Shadow Configuration Drift occurs when source code silently overrides the intent of the
higher-level configuration. Given a dependency relation $D \subseteq C \times C$, it is the set
of inconsistencies that can be detected only through the transitive closure $D^*$:
\begin{equation}
\mathit{ShadowDrift} = \Big\{ d \in \mathit{ConfigDrift} : d \notin \bigcup_i \mathit{SingleFile}(C_i)
\wedge d \in \mathit{CrossFile}(D^*) \Big\}
\end{equation}

\subsection{Motivating Examples}

\begin{description}
\item[QoS override.] The YAML file sets \code{reliability: RELIABLE} for a sensor topic, but the
  C++ node creates it with \code{rclcpp::QoS(10).best\_effort()}. Under network load, messages
  are dropped without any warning.
\item[Security bypass.] The XML governance file requires \code{ENCRYPT} protection, but a Python
  launch script binds the node to \code{127.0.0.1}, which bypasses the network-level policy.
\item[Lifecycle violation.] The launch file sets up a managed lifecycle for a motor controller,
  but the C++ code publishes commands before the node has gone through its state transitions.
\item[Ignored QoS override.] A parameter file sets
  \code{qos\_overrides./scan.subscription.reliability: best\_effort}, but the subscription was
  created without \code{QosOverridingOptions}. The override is ignored and the QoS in the code
  applies. Finding this needs three inputs: the parameter file, the launch graph that binds it
  to the node, and the node's source.
\item[Namespaced parameter keys.] \code{rcl} matches parameter-file keys against fully
  qualified node names. A section keyed \code{amcl} therefore does not apply to an AMCL node
  launched in namespace \code{robot1} (\code{/robot1/amcl}), and its values are silently lost.
\end{description}

% ===========================================================================
\section{Design of rosfix}
\label{sec:design}

\subsection{Pipeline}

rosfix analyzes a workspace in four stages:
\begin{enumerate}
\item \textbf{Parsing.} Each language has its own parser: roxmltree for XML, serde\_yaml for
  YAML, rustpython-parser for Python and Tree-sitter for C++.
\item \textbf{Intermediate representation.} The parsed files are mapped to one representation
  of ROS 2 entities (nodes, parameters, topics, QoS).
\item \textbf{Dependency resolution.} A cross-file dependency graph is built, and parameter
  values are tracked as they propagate through it.
\item \textbf{Drift detection.} Patterns are matched across the configuration sources.
\end{enumerate}

\subsection{Cross-Language Intermediate Representation}

The intermediate representation holds ROS 2 entities independently of the file format they came
from:
\begin{lstlisting}
struct CrossLanguageIR {
  nodes: Vec<NodeDefinition>,
  parameters: Vec<ParameterDefinition>,
  topics: Vec<TopicDefinition>,
  qos_profiles: Vec<QoSProfile>,
  dependencies: DependencyGraph,
  file_origins: HashMap<ElementId, FileOrigin>
}
\end{lstlisting}
Every element records the file and position it came from, so a finding can point to the exact
location in each file involved.

\subsection{Drift Patterns}

rosfix detects seven drift patterns:
\begin{description}
\item[QoS override:] QoS hardcoded in source code that contradicts the configuration files.
\item[Security policy bypass:] source code that circumvents the SROS2 security configuration.
\item[Parameter conflict:] different values for one parameter in different layers.
\item[Lifecycle violation:] code that does not follow the managed lifecycle it is launched with.
\item[Topic type mismatch:] publishers and subscribers with incompatible message types.
\item[Executor safety:] blocking calls inside callbacks of a single-threaded executor.
\item[Physical safety:] URDF joint limits and navigation footprints.
\end{description}

\subsection{Fact Extraction}
\label{sec:facts}

Each file is reduced to typed facts. rosfix has extractors for:
\begin{itemize}
\item \code{package.xml}, \code{setup.py} and \code{setup.cfg};
\item CMake, through a small interpreter for variables, executables and libraries,
  \code{install} rules, ament\_cmake\_auto, rclcpp component registration and
  generate\_parameter\_library specifications;
\item launch files in XML, YAML and Python, and parameter YAML files;
\item C++ through Tree-sitter~\cite{treesitter}, and Python through the RustPython AST;
\item URDF, with a subset of xacro (properties, arguments, arithmetic, macros with block
  parameters, includes and conditionals);
\item SROS2 governance, permission and policy documents, and RViz configurations, whose
  displays set their own QoS.
\end{itemize}

Python launch files are not executed. A static interpreter evaluates the common subset
(\code{Node}, \code{IncludeLaunchDescription}, \code{GroupAction}, \code{PushRosNamespace},
\code{SetParameter}, \code{DeclareLaunchArgument}, \code{LaunchConfiguration}, path
substitutions, helper functions, \code{RewrittenYaml} and \code{ParameterFile}) and marks
everything else as unknown. The YAML layer keeps exact byte positions and types scalars the way
\code{rcl\_yaml\_param\_parser} does: \code{1} is an integer, \code{1.0} a double and \code{on}
a boolean, as at runtime. From node code, rosfix records node names, declared parameters with
their types and defaults, endpoints with literal names and statically evaluated QoS, callbacks,
callback groups, executors and blocking calls.

\subsection{Resolved Launch Graph}
\label{sec:graph}

Every launch file is the root of a graph. The resolver follows includes into the workspace and
into installed packages. It propagates launch arguments (a value passed by the including file
overrides the default), composes namespace pushes (relative pushes nest, absolute ones reset),
and applies remappings to the fully expanded names. Parameter-file sections are matched as in
\code{rcl}: a key $k$ applies to node $n$ if and only if
\begin{equation}
\mathrm{match}\big(\mathrm{abs}(k), \mathrm{fqn}(n)\big),\quad
\mathrm{abs}(k) = \begin{cases} k & k \text{ starts with } / \\ /k & \text{otherwise}\end{cases}
\end{equation}
where \code{*} matches one name token and \code{**} any number of tokens. For
\code{RewrittenYaml(root\_key=$r$)}, keys are prefixed with $r$. Each launched node is linked to
its implementation through the package's build rules: the sources of the executable, in-package
libraries linked into it, and node classes defined in headers. Composable nodes are linked
through their plugin class.

\subsection{Precision First}
\label{sec:precision}

Every fact records what could not be determined, and every rule checks these markers before it
reports. Parameter declarations are in one of three states:
\begin{itemize}
\item \emph{closed}: all declared names are known;
\item \emph{dotted-open}: names that contain a dot are uncertain, as with plugin parameters
  such as \code{plugin\_name.param};
\item \emph{open}: at least one name is fully dynamic.
\end{itemize}
Declarations made through another object's node handle are attributed to a node class when
possible. Literal declarations in the node's workspace dependencies are accepted. A node that
derives from a base class rosfix cannot find is treated as open. Launch graphs with unresolved
includes or dynamic code are marked partial, and checks that need complete knowledge are skipped
for them.

\subsection{Installed Packages and Knowledge Packs}
\label{sec:installed}

When ROS is sourced, rosfix reads the ament resource index~\cite{ament_index} of every prefix in
\code{AMENT\_PREFIX\_PATH}. This gives the installed packages, their executables
(\code{lib/<pkg>}), pluginlib plugins (\code{<category>\_\_pluginlib\_\_plugin}), rclcpp
components and dependencies. Launch and parameter files of installed packages that the workspace
includes are loaded for resolution, but no finding is reported inside them.

Binary installs contain no source code, so the parameters of installed nodes are unknown. rosfix
therefore embeds a \emph{knowledge pack} per distribution, built by running its own extractors
once over the source of commonly used packages (Table~\ref{tab:packs}). A pack is used only when
the distribution is known, because parameter names change between distributions.

A pack also records each node's topic endpoints with the QoS its source requests. An endpoint
created with \code{SystemDefaultsQoS} or \code{rmw\_qos\_profile\_system\_default} is stored as
\emph{system default}, not as the ROS default profile: its reliability and durability are chosen
by the middleware (and by a DDS profile file, if one is loaded), so the QoS agent treats them as
unknown unless the launch environment names such a file. Each distribution also embeds the list
of keys \code{rosdep} can resolve: the packages released in that distribution plus rosdep's
system keys (4{,}590 for Humble, 4{,}696 for Jazzy, 3{,}990 for Kilted), which rule L009 checks
against.

\begin{table}[t]
\centering
\caption{Embedded knowledge packs, built from nine upstream repositories at each distribution's
branch (rebuilt 24 November 2026). Endpoints: topic endpoints with QoS, of which
\emph{sys.} request the system-default profile.}
\label{tab:packs}
\footnotesize
\setlength{\tabcolsep}{4pt}
\begin{tabular}{@{}lrrrrrrr@{}}
\toprule
Distro & Pkgs & Exec. & Comp. & Params & Endpoints & Sys. & KiB\\
\midrule
Humble & 42 & 100 & 54 & 2{,}335 & 290 & 4 & 14\\
Jazzy  & 44 &  96 & 38 & 3{,}051 & 240 & 6 & 16\\
Kilted & 47 &  97 & 40 & 2{,}985 & 232 & 8 & 16\\
\bottomrule
\end{tabular}
\end{table}

\subsection{Rules}
\label{sec:rules}

Domain agents run over the facts and the launch graphs (Section~\ref{sec:incremental}). At the
time of writing there are 18 agents (17 domain agents and \code{core}, which reports files that
could not be parsed) and 93 rules, each with a stable identifier (Table~\ref{tab:rules};
\code{rosfix rules} prints the current list). A finding carries an exact location, related
locations, an explanation and, where possible, a fix. Rules T005 and S003 are off by default.

\begin{table*}[t]
\centering
\caption{Agents and rules (R001 reports files that could not be parsed)}
\label{tab:rules}
\footnotesize
\begin{tabularx}{\textwidth}{@{}l l X@{}}
\toprule
Agent & Rules & What is checked\\
\midrule
params    & P001--P008 & parameter-file key matches no node; missing \code{ros\_\_parameters}; undeclared parameter; type mismatch; required parameter not set; value outside a declared range or validator; parameter renamed or removed between ROS distributions; two parameters of one node that contradict each other\\
topics    & T001--T005 & near-miss topic names; remap without effect; message type conflict; absolute name inside a namespace; topic without publisher\\
executor  & E001--E002 & blocking wait or spin in a callback that shares the response's callback group or a single-threaded executor; sleep in a callback\\
control   & C001--C010 & unknown joint; missing command or state interface; unknown controller type; inconsistent controller configuration; \code{<ros2\_control>} joint missing from URDF; missing hardware plugin; spawned controller not defined; defined controller never spawned; value violating the controller's parameter specification; controller \code{update\_rate} the manager cannot run\\
packaging & L001--L012 & launched executable or component not built, registered or installed; undeclared dependency; wrong script directory in \code{setup.cfg}; uninstalled files used at runtime; missing file; component class without \code{RCLCPP\_COMPONENTS\_REGISTER\_NODE}; interface files not generated; missing \code{rosidl\_interface\_packages} group; dependency that \code{rosdep} cannot resolve; component library never installed; \code{install()} rule naming a path that does not exist; Python subpackage left out of a hand-written \code{setup.py}\\
simtime   & S001--S003 & mixed \code{use\_sim\_time}; simulator launched without sim time; sim time without a clock source\\
qos       & Q001--Q005 & incompatible reliability, durability, deadline or liveliness; late-joining subscriber; ignored \code{qos\_overrides}\\
nav2      & N001--N006, N008--N010 & plugin section missing; unknown plugin type; unknown frame; footprint; inflation radius too small; unknown parameter; plugin renamed between distributions; unknown behavior-tree node type; unknown behavior-tree subtree (there is no N007)\\
lifecycle & M001--M002 & node not managed by a lifecycle manager; lifecycle manager lists an unknown node\\
tf        & F001--F005 & frame with two parents; unparsable \code{static\_transform\_publisher} arguments; lookup frame with a leading slash; unknown frame; cycle in the static transform tree\\
launch    & A001--A003 & launch configuration used but never declared or passed; declared argument never used; two nodes of one launch graph with the same fully qualified name\\
network   & D001--D005 & invalid \code{ROS\_DOMAIN\_ID}; unknown \code{RMW\_IMPLEMENTATION}; missing or invalid DDS configuration file; deprecated \code{ROS\_LOCALHOST\_ONLY}; invalid \code{rmw\_zenoh\_cpp} environment value\\
security  & X001--X006 & weak governance; broad or expired permissions; unknown enclave; policy does not match the launch graph; absolute keystore path\\
moveit    & V001--V005 & \code{kinematics.yaml} group not in the SRDF; joint limit for an unknown joint; integer joint limit; SRDF names an unknown joint or link; controller not defined by any ros2\_control configuration\\
urdf      & U001--U004 & element that urdfdom rejects; physically impossible mass or inertia; missing \code{package://} mesh (or one that differs only in case); empty joint-limit range\\
gzbridge  & G001--G005 & ROS/Gazebo type pair the bridge cannot convert; bridge YAML entry rejected at start; malformed \code{parameter\_bridge} argument; unknown YAML key; one-way bridge pointing against the ROS nodes' use\\
ros1      & O001--O003 & ROS 1 XML launch file, ROS 1 client-library code, or catkin package inside a ROS 2 workspace\\
\bottomrule
\end{tabularx}
\end{table*}

\subsection{Demand-Driven, Incremental Analysis}
\label{sec:incremental}

Agents are not invoked in a fixed sequence over everything. Each agent declares (i) the file
kinds its findings depend on, directly or through the launch graph, and (ii) a cheap trigger
predicate over the facts (the control agent has no work without a \code{<ros2\_control>} block
or a controller-manager configuration). A re-analysis after an edit (fix verification, the
mutation benchmark, \code{rosfix watch}, the language server) is treated as a \emph{change
event}: the engine compares per-file content hashes with the previous run, re-parses only the
edited files (a parse cache keyed by file and content hash; URDF/xacro is never cached because
it includes other files), rebuilds only the launch graphs whose launch or parameter files
changed, and runs only the agents whose declared inputs intersect the changed kinds. The others
reuse their previous findings. Graph reuse is disabled when code, manifests or build files
change, because node-to-source linking depends on them.

The scheme is only as sound as the declared inputs. A differential test edits every file of a
sandbox workspace in several ways (79 edits) and requires the incremental result to equal a
fresh full analysis; removing a single declared input (URDF from the control agent) or not
recording the parameter files a launch graph reads makes the test fail.

\subsection{Log Triage}
\label{sec:triage}

Static analysis runs before a launch; developers usually meet the same faults afterwards, as an
error in a log. \code{rosfix triage <log>} recognises the runtime messages that the faults of
Section~\ref{sec:rules} produce, using the message formats of pluginlib, rclcpp\_components,
launch, the Nav2 lifecycle manager, rosdep, rclcpp's QoS events, BehaviorTree.CPP and the
controller manager, and links each recognised line to the findings that explain it (rule,
location, message). For example, ``class \code{differential} with base class type
\code{nav2\_amcl::MotionModel} does not exist'' is explained by N008 on the AMCL
\code{robot\_model\_type} line that still uses the pre-Humble name.

\subsection{Verified Repair}
\label{sec:repair}

An agent attaches a fix only when the edit is local and mechanical, for example \code{1} to
\code{1.0}, renaming a misspelled key to the closest declared name, or adding a
\code{<depend>}. Fixes are applied in memory and the whole workspace is analyzed again. A fix is
kept only if its finding disappears and no new error appears in a file it touched
(Algorithm~\ref{alg:fix}). A file that no longer parses counts as a new error, even though a
parse failure is otherwise reported only at information level; Section~\ref{sec:fixsandbox} shows
why this matters. Kept fixes are written atomically, and only when the user passes
\code{--write}.

\begin{algorithm}[t]
\caption{Fix verification}
\label{alg:fix}
\footnotesize
\begin{algorithmic}[1]
\State $F \gets$ non-overlapping fixes of all findings $D$
\For{round $=1..3$}
  \State $D' \gets \mathrm{analyse}(\text{workspace with } F \text{ applied in memory})$
  \State $R \gets \{f\in F \mid \mathrm{finding}(f)\in D' \lor \mathrm{newErrors}(D', D, f)\}$
  \If{$R=\emptyset$} \Return $F$ \EndIf
  \State $F \gets F \setminus R$
\EndFor
\State \Return $F$
\end{algorithmic}
\end{algorithm}

\textbf{Repair assistant.} Some findings have no mechanical fix, for example a key
\code{max\_lin\_vel} for a node that declares \code{max\_linear\_velocity}. These findings carry
a repair site: the exact text span and the set of values that are valid there (declared
parameters, \code{<ros2\_control>} joints, URDF frames, executables or enclaves). The assistant
picks one of these values or abstains. It never generates text, unlike generative repair
models~\cite{silva2023repairllama,wang2023rapgen}. Two rankers score the candidates: a lexical
ranker (edit distance, word overlap and common ROS abbreviations) and a learned ranker, a
16-16-1 perceptron over 16 string features with 289 weights compiled into the binary. A choice
needs a clear margin over the next candidate and must pass Algorithm~\ref{alg:fix}. It is written
only with an explicit flag, because re-analysis cannot confirm what the user intended. We first
tried a code language model with 0.5B parameters on the CPU; it needed 4.5\,s to load and
17.6\,s per finding, which is why rosfix uses the small ranker instead.

\textbf{Sandboxed build check.} With \code{--verify-build}, rosfix copies the workspace twice
into a temporary directory and builds both copies with \code{colcon}, one before and one after
the fixes. When Docker is available, the builds run in a ROS container without network access.
Fixes that break a build that worked before are refused.

\textbf{Optional references.} With \code{--online}, findings get links to related Robotics Stack
Exchange questions and GitHub issues from the main ROS organizations. Queries are built only from
the rule that fired, so no file content, path or name leaves the machine. A separate helper
process sends the requests to an HTTPS allowlist with timeouts and size limits, and the results
are cached.

\subsection{Implementation and Distribution}

rosfix is written in Rust (about 16{,}700 lines) and runs single-threaded by default. It reports
text, JSON and SARIF~2.1.0~\cite{sarif} with exact regions, rule metadata and fixes. It supports
inline suppressions and baselines for gradual adoption, and it exits with a non-zero code at a
configurable severity. It is distributed as Python wheels (with a \code{ros2 rosfix} command),
standalone archives, container images, a GitHub Action and pre-commit hooks.

The command line is organised as subcommands (\code{check}, \code{fix}, \code{graph},
\code{watch}, \code{lsp}, \code{baseline}, \code{init}, \code{doctor}, \code{migrate}, \code{triage},
\code{docs}); the original single-command form still works unchanged. Two commands reuse the
analysis in new ways. \code{migrate --to} analyzes a workspace with the knowledge packs of its
current and of a target distribution and reports only the findings that are new on the target,
such as parameters that upstream renamed between releases. \code{doctor} checks the middleware
environment of the user's shell (\code{RMW\_IMPLEMENTATION}, \code{ROS\_DOMAIN\_ID}, discovery
settings, DDS profile files), which depends on the machine and is therefore kept out of
\code{check}. The manual is compiled into the binary (\code{rosfix docs}); its rule reference is
generated from the rule registry, so it cannot drift from the implementation.

% ===========================================================================
\section{Evaluation}
\label{sec:eval}

Sections~\ref{sec:dataset} to~\ref{sec:perf} report measurements of an earlier version of
rosfix. Sections~\ref{sec:current} to~\ref{sec:coverage} were measured on the current
implementation described in Section~\ref{sec:design}.

\subsection{Dataset and Setup}
\label{sec:dataset}

We evaluated rosfix on 12 open-source ROS 2 repositories with 38,917 files in total
(Table~\ref{tab:dataset}). All experiments ran on Windows (build 10.0.26200) on an AMD Ryzen
machine, with release builds made by Rust 1.97.1.

\begin{table}[t]
\centering
\caption{Repositories in the evaluation}
\label{tab:dataset}
\footnotesize
% CHECK: the per-repository counts below sum to 39,117, not 38,917. Please verify before submission.
\begin{tabular}{@{}lr@{\hspace{2em}}lr@{}}
\toprule
Repository & Files & Repository & Files\\
\midrule
PX4-Autopilot  & 26{,}404 & rmw\_fastrtps   & 224\\
Autoware       &  8{,}074 & demos           & 217\\
MoveIt2        &  1{,}567 & system\_tests   & 104\\
Navigation2    &  1{,}333 & turtlebot3      &  88\\
rmf\_ros2      &     695  & sros2           &  73\\
ros2\_control  &     314  & rmw\_cyclonedds &  24\\
\midrule
\multicolumn{3}{@{}l}{Total} & 38{,}917\\
\bottomrule
\end{tabular}
\end{table}

\subsection{Detection Accuracy}
\label{sec:accuracy}

rosfix reported 2,405 configuration issues in 947 files. We validated a stratified random sample
of 312 issues by hand. Table~\ref{tab:accuracy} gives the precision with its 95\% confidence
interval and the issues by severity.

\begin{table}[t]
\centering
\caption{Detection results of the earlier version}
\label{tab:accuracy}
\footnotesize
\begin{tabular}{@{}lr@{}}
\toprule
Measure & Value\\
\midrule
Issues reported & 2{,}405\\
Files with issues & 947\\
Precision ($n=312$, 95\% CI) & 87.3\% (83.1--91.5\%)\\
\midrule
Critical (security, physical safety) & 127\\
Error (functional problems) & 891\\
Warning (performance, consistency) & 1{,}387\\
\bottomrule
\end{tabular}
\end{table}

\subsection{Shadow Drift}
\label{sec:drift}

Of the 2,405 issues, 312 (13.0\%) are shadow drift: they appear only when several files are
analyzed together (Table~\ref{tab:drift}). Single-file analyzers found none of them.

\begin{table}[t]
\centering
\caption{Shadow drift instances by pattern}
\label{tab:drift}
\footnotesize
\begin{tabular}{@{}lrl@{}}
\toprule
Pattern & Cases & Typical form\\
\midrule
QoS override        & 189 & reliability downgraded in code\\
Parameter conflict  &  76 & competing values across layers\\
Security bypass     &  47 & policy circumvented\\
\midrule
Total               & 312 & \\
\bottomrule
\end{tabular}
\end{table}

\subsection{Performance}
\label{sec:perf}

Table~\ref{tab:scal} lists the analysis time per file. With Rayon, the analysis ran
3.19$\times$ faster on 4 cores and 5.15$\times$ faster on 8 cores, a parallel efficiency of 80\%
and 64\%.

\begin{table}[t]
\centering
\caption{Analysis time of the earlier version}
\label{tab:scal}
\footnotesize
\begin{tabular}{@{}llrr@{}}
\toprule
Size & Repository & Files & ms/file\\
\midrule
Small      & demos            &    217 &  9.55\\
Large      & MoveIt2          &  1{,}567 & 21.94\\
Very large & Autoware         &  8{,}074 & 20.06\\
\midrule
All        & 12 repositories  & 38{,}917 & 13.81\\
\bottomrule
\end{tabular}
\end{table}

\subsection{Current Implementation on the 12 Repositories}
\label{sec:current}

We ran the current implementation (release build, one thread, no ROS installation) on shallow
clones of the same 12 repositories. Table~\ref{tab:current} gives the median of three runs after
one warm-up run and the peak memory. The machine was shared with other work during the runs, so
the times are indicative; repeated runs differed by up to a factor of two. The file counts are
lower than in Table~\ref{tab:dataset} because rosfix now parses only files that belong to a ROS
package (a directory with a \code{package.xml}) and that pass a check of type and content:
manifests, build files, launch, parameter, URDF and source files. The Autoware
repository is a meta-repository whose packages are fetched separately, so the clone contains no
ROS package.

\begin{table}[t]
\centering
\caption{Current implementation: one thread, median of three runs}
\label{tab:current}
\footnotesize
\begin{tabular}{@{}lrrrl@{}}
\toprule
Repository & Files & Time (ms) & Memory (MB) & Findings\\
\midrule
navigation2    & 818 & 2{,}139 & 65 & L002$\times$34, S002$\times$10\\
moveit2        & 492 & 1{,}659 & 28 & L002$\times$32\\
demos          & 145 &     243 & 13 & L002$\times$6, E002$\times$1\\
rmf\_ros2      & 140 &     542 & 17 & L002$\times$2, P003$\times$3\\
ros2\_control  & 117 &     631 & 20 & L002$\times$1\\
system\_tests  &  84 &     150 & 12 & none\\
turtlebot3     &  63 &     101 & 12 & L002$\times$5, E002$\times$2\\
PX4-Autopilot  &  25 &     634 & 13 & none\\
sros2          &  20 &      67 &  7 & none\\
rmw\_fastrtps  &  17 &      61 & 10 & none\\
rmw\_cyclonedds &  1 &     130 & 14 & none\\
Autoware       &   0 &      24 &  6 & none\\
\bottomrule
\end{tabular}
\end{table}

We reviewed all 96 findings by hand (Table~\ref{tab:precision}) and found no false positive. All 80 L002 findings are
dependencies that a package uses but does not declare in \code{package.xml}. For example,
\code{turtlebot3\_bringup} starts nodes from \code{rclcpp\_components}, \code{camera\_ros} and
\code{image\_view} without declaring them, so installing the package from binaries does not pull
them in. The S002 findings are Nav2 example nodes started next to a simulator without
\code{use\_sim\_time}. They match the rule's definition, but whether they cause a failure depends
on how the nodes use time. The three E002 findings are sleeps inside callbacks of
single-threaded executors. The three P003 findings are parameters (\code{subns}, \code{verbose}
and \code{run\_composition\_command}) that the rmf\_ros2 launch file passes to
\code{robot\_state\_aggregator} but that the node never declares.

\begin{table}[t]
\centering
\caption{Manual review of the current findings}
\label{tab:precision}
\footnotesize
\begin{tabular}{@{}llrr@{}}
\toprule
Rule & Meaning & Findings & Correct\\
\midrule
L002 & undeclared dependency       & 80 & 80\\
S002 & simulator without sim time  & 10 & 10\\
E002 & sleep in callback           &  3 &  3\\
P003 & undeclared parameter        &  3 &  3\\
\midrule
Total &                            & 96 & 96\\
\bottomrule
\end{tabular}
\end{table}

Running on real code exposed false-positive sources that the unit tests had missed. Each was
fixed in the extractors, the resolver or the agents:
\begin{itemize}
\item parameter files passed through \code{ParameterFile(RewrittenYaml(\dots))};
\item nested dictionaries in Python launch parameters, which ROS flattens to dotted names;
\item project-specific QoS classes such as \code{LatchedPublisherQoS};
\item launch\_testing files, and ROS 1 \code{.launch} and \code{.test} files left in test
  directories;
\item callbacks that hand their work to a thread, and sleeps in action callbacks of a
  multi-threaded rclpy executor;
\item metapackage dependencies, and parameters declared by shared library code in dependency
  packages;
\item the member form of \code{declare\_or\_get\_parameter}, and template calls that are
  misparsed without type information;
\item frame ids with a leading slash on the publishing side, which tf2 strips on receipt;
\item parameters declared through generate\_parameter\_library under a namespace passed to
  \code{ParamListener} at runtime, as in the MoveIt Servo demos.
\end{itemize}

\subsection{Mutation Benchmark}
\label{sec:mutation}

The benchmark injects one fault at a time, in memory, at a site found in the clean analysis. A
site is real code: for example, a parameter that a node declares and a launch file binds. A
mutation counts as detected only if the expected rule fires at the mutated line, or, for a
removed dependency, names that dependency. At most 20 sites per operator and repository were
used. We re-ran the benchmark on 24 November 2026 with the current rule set and packs, on
shallow clones of the 12 repositories at their default branches (the revisions are listed in
\code{paper/data/mutation\_bench\_revisions.txt}). All 161 mutations were detected
(Table~\ref{tab:mut}); an earlier run on older revisions had found 165 sites, all detected. The
site counts change with the code of the repositories, not with rosfix. We also injected three faults by
hand into Nav2's default parameter file: a misspelled \code{controller\_frequency}, an integer
where a double is expected, and a misspelled \code{ros\_\_parameters}. rosfix detected each one
and produced a verified fix for it.

\begin{table}[t]
\centering
\caption{Mutation benchmark on the 12 repositories}
\label{tab:mut}
\footnotesize
\begin{tabular}{@{}llrr@{}}
\toprule
Mutation & Rule & Injected & Detected\\
\midrule
\code{ros\_\_parameters} misspelled  & P002 &  50 &  50\\
Nav2 plugin type misspelled          & N002 &  40 &  40\\
executable name misspelled in launch & L001 &  28 &  28\\
dependency removed from package.xml  & L002 &  16 &  16\\
parameter key misspelled             & P003 &  20 &  20\\
double written as integer            & P004 &   4 &   4\\
parameter key abbreviated            & P003 &   3 &   3\\
\midrule
Total                                &      & 161 & 161\\
\bottomrule
\end{tabular}
\end{table}

\subsection{Repair Assistant}
\label{sec:assistant}

The learned ranker is trained on synthetic corruptions of the real parameter names and
executables in the knowledge packs (Table~\ref{tab:packs}): typos, double typos, abbreviations,
abbreviation plus typo, dropped or swapped words, camelCase and plurals. A further set of
\emph{answer-absent} cases removes the intended name from the candidates; any pick there is a
wrong edit, so these cases teach abstention. Cases are assigned to five folds by package. A run
holds out one fold for testing, uses the next for calibration and trains on the other three, and
a held-out case is dropped if its name also occurs in training: packages share many names (for
example the parameters of Nav2's base classes), so a split by package alone lets the ranker see
most test names during training. Training takes 40 epochs of Adam on a listwise loss;
calibration then lowers all scores until the learned ranker makes no more wrong edits than the
lexical one on the calibration fold. Fold~0 is the shipped model. \code{rosfix-train} writes the
weights only if, summed over all five folds, the learned ranker is right more often than the
lexical ranker \emph{and} makes no more wrong edits. Training and evaluation run in plain Rust on
a CPU and are deterministic for a given seed (two runs wrote identical weights).

Table~\ref{tab:assistheld} gives the shipped model (fold~0; 5,619 training and 1,370 calibration
cases) on its 1,791 name-disjoint held-out cases, with the packs rebuilt on 24 November 2026;
4,671 further held-out cases shared their name with training and were excluded. The learned
ranker chose the correct value in 738 of 1,062 cases where it was offered, against 575 for the
lexical ranker (69.5\% vs.\ 54.1\%), and made 7 wrong edits in all 1,791 cases, against 30
(0.39\% vs.\ 1.68\%). On answer-absent cases it abstained 724 times out of 729 (99.3\%), against
704 (96.6\%). The lexical ranker is better on dropped words (63 vs.\ 9 correct), where the
learned ranker abstains instead, and slightly better on camelCase (105 vs.\ 102).

The fold-0 result depends on which packages fall into the test fold. Table~\ref{tab:assistcv}
repeats training and evaluation for all five folds. Pooled over the 3,005 name-disjoint held-out
cases, the learned ranker was correct 1,298 times against 801 (43.2\% vs.\ 26.7\%) and made 48
wrong edits against 59 (1.60\% vs.\ 1.96\%). It is correct more often in every fold, but in folds
1, 2 and 4 it makes \emph{more} wrong edits than the lexical ranker; its overall advantage on
wrong edits comes from fold~0, the largest. Folds 2 and 3 are small (140 and 38 cases), so their
rates are noisy. Reducing wrong edits on unseen packages, rather than raising accuracy, is
therefore the open problem for the assistant; every edit is still verified by re-analysis and
written only on request.

\begin{table}[t]
\centering
\caption{Shipped repair ranker (fold 0) on name-disjoint held-out cases (correct / abstained /
wrong). For answer-absent cases abstaining is correct and any pick is wrong.}
\label{tab:assistheld}
\footnotesize
\setlength{\tabcolsep}{4pt}
\begin{tabular}{@{}lrcc@{}}
\toprule
Corruption & $n$ & Lexical & Learned\\
\midrule
typo              & 143 &  69\,/\,74\,/\,0  & \textbf{133}\,/\,10\,/\,0\\
double typo       & 139 &  16\,/\,123\,/\,0 & \textbf{78}\,/\,61\,/\,0\\
abbreviation      & 148 &  67\,/\,80\,/\,1  & \textbf{110}\,/\,38\,/\,0\\
abbrev.\ + typo   & 138 &  16\,/\,121\,/\,1 & \textbf{34}\,/\,102\,/\,2\\
plural            & 162 & 126\,/\,36\,/\,0  & \textbf{151}\,/\,11\,/\,0\\
swapped words     & 127 & 113\,/\,14\,/\,0  & \textbf{121}\,/\,6\,/\,0\\
camelCase         & 106 & \textbf{105}\,/\,1\,/\,0 & 102\,/\,4\,/\,0\\
dropped word      &  99 & \textbf{63}\,/\,33\,/\,3 & 9\,/\,90\,/\,0\\
\midrule
answer offered    & 1{,}062 & 575\,/\,482\,/\,5 & \textbf{738}\,/\,322\,/\,\textbf{2}\\
absent: siblings  &   546 & --\,/\,523\,/\,23 & --\,/\,\textbf{542}\,/\,4\\
absent: unrelated &   183 & --\,/\,181\,/\,2  & --\,/\,\textbf{182}\,/\,1\\
\midrule
all               & 1{,}791 & 575\,/\,1{,}186\,/\,30 & \textbf{738}\,/\,1{,}046\,/\,\textbf{7}\\
\bottomrule
\end{tabular}
\end{table}

\begin{table}[t]
\centering
\caption{Five-fold cross-validation over package folds, name-disjoint held-out cases: share of
held-out cases (answer offered or absent) with a correct choice, and with a wrong edit}
\label{tab:assistcv}
\footnotesize
\begin{tabular}{@{}lrrrrr@{}}
\toprule
 & & \multicolumn{2}{c}{Correct (\%)} & \multicolumn{2}{c}{Wrong edits (\%)}\\
\cmidrule(lr){3-4}\cmidrule(l){5-6}
Fold & $n$ & Lexical & Learned & Lexical & Learned\\
\midrule
0 & 1{,}791 & 32.1 & \textbf{41.2} & 1.68 & \textbf{0.39}\\
1 &     730 & 14.7 & \textbf{39.0} & \textbf{2.05} & 2.88\\
2 &     140 & 25.7 & \textbf{52.9} & \textbf{2.14} & 5.71\\
3 &      38 & 31.6 & \textbf{57.9} & 5.26 & 5.26\\
4 &     306 & 23.2 & \textbf{58.5} & \textbf{2.94} & 3.27\\
\midrule
mean $\pm$ sd & & 25.5\,$\pm$\,7.1 & \textbf{49.9}\,$\pm$\,9.2 & \textbf{2.8}\,$\pm$\,1.4 & 3.5\,$\pm$\,2.1\\
pooled & 3{,}005 & 26.7 & \textbf{43.2} & 1.96 & \textbf{1.60}\\
\bottomrule
\end{tabular}
\end{table}

Table~\ref{tab:assist} shows both rankers on the benchmark mutations of
Section~\ref{sec:mutation} whose correct value is known. The learned ranker fixed all 28
misspelled executable names, against 12 for the lexical ranker. On misspelled parameter keys both
rankers abstained in all 20 cases, although the correct key was among the candidates in 19 of
them. Abstaining leaves the finding reported, so it is safe, but this is the main gap in the
assistant.

\begin{table}[t]
\centering
\caption{Repair choices on benchmark mutations (correct / abstained / wrong)}
\label{tab:assist}
\footnotesize
\begin{tabular}{@{}lrrcc@{}}
\toprule
Mutation & Cases & In set & Lexical & Learned\\
\midrule
executable name misspelled & 28 & 28 & 12\,/\,16\,/\,0 & 28\,/\,0\,/\,0\\
parameter key misspelled   & 20 & 19 &  0\,/\,20\,/\,0 &  0\,/\,20\,/\,0\\
parameter key abbreviated  &  3 &  3 &  3\,/\,0\,/\,0  &  3\,/\,0\,/\,0\\
\bottomrule
\end{tabular}
\end{table}

\subsection{Fix Verification Sandbox}
\label{sec:fixsandbox}

Algorithm~\ref{alg:fix} checks that a finding disappears; it does not check that the written
file is still what the user expects. We therefore test \code{rosfix fix} end to end, on the
binary, in two settings.

\textbf{Planted workspace.} A test workspace contains one mechanically fixable mistake for each
of the 12 rules that have safe fixes (P001--P004, L002--L004, N002, N003, N006, C003, F003), 17
fixes in 11 files, plus traps that must stay untouched: a finding hidden by a suppression
comment, an ambiguous misspelling with two close candidates, a one-line \code{package.xml}, and
a \code{package.xml} with Windows line endings. Eleven tests check that a dry run writes
nothing; that \code{--write} produces exactly the hand-reviewed file contents, byte for byte,
and changes no other file; that the printed diff equals the change on disk; that every fixed
finding is gone and no new finding appears; that a second run finds nothing to fix; that the
old flag form and the \code{fix} command write identical files; that with a machine-readable
report the diff goes to stderr; that the exit code counts only findings left after writing; and
that findings hidden by a baseline or needing the assistant are not written.

Building this workspace exposed four defects, all fixed. (1) In a \code{package.xml} written on
one line, an added dependency was inserted after \code{</package>}, outside the root element.
Verification accepted it, because a file that does not parse is reported only at information
level; the fix itself had made the dependency finding disappear. Fixes now insert inside the
element, and a new parse failure in a touched file rejects the fix. (2) A parameter file whose
only node section misspells \code{ros\_\_parameters} was not recognised as a parameter file, so
P002 never fired; because the file was then unread, the loading node's parameter sources counted
as incomplete, which silently turned off P001 for every file that node loads. (3) A trailing
suppression comment also hid findings on the following line; it now applies to its own line
only, while a comment on a line of its own still covers the next line. (4) A dependency added by
an L002 fix can itself be reported, for example as a package that \code{rosdep} cannot resolve;
a fix is now also rejected when any new finding, of any severity, points at the text it wrote.

\textbf{Real repositories.} We ran \code{rosfix fix --write} on copies of the 12 repositories of
Table~\ref{tab:dataset} with the Jazzy pack and checked the same properties on each.
Table~\ref{tab:fixcorpus} shows the result. Of 141 findings, 90 had a verified safe fix, written
to 41 files; one candidate fix was rejected by the rule in (4). After writing, no new finding
appeared, every changed XML file parsed, every printed diff equalled the change on disk, and a
second run found nothing to fix. All 90 written fixes add a missing dependency declaration
(L002); the other fixable rules had no occurrences in these repositories. The counts are higher
than in Table~\ref{tab:current} because this run uses the Jazzy pack and the current rule set.
These checks establish that each written file is well formed and that the edit does what the
finding asks; they do not establish that every added dependency is what the maintainers would
have written, which the hand review in Section~\ref{sec:current} covers for the 80 L002 findings
reviewed there. A dry run took at most 2.2\,s per repository.

\begin{table}[t]
\centering
\caption{\code{rosfix fix --write} on the 12 repositories (Jazzy pack)}
\label{tab:fixcorpus}
\footnotesize
\begin{tabular}{@{}lrrrr@{}}
\toprule
Repository & Findings & Fixed & Files & Rejected\\
\midrule
navigation2   & 53 & 40 & 12 & 0\\
moveit2       & 40 & 31 & 16 & 1\\
turtlebot3    & 14 &  8 &  3 & 0\\
demos         &  9 &  8 &  7 & 0\\
rmf\_ros2     & 10 &  2 &  2 & 0\\
ros2\_control &  1 &  1 &  1 & 0\\
other six     & 14 &  0 &  0 & 0\\
\midrule
total         & 141 & 90 & 41 & 1\\
\bottomrule
\end{tabular}
\end{table}

\subsection{Validation Protocol for New Checks}
\label{sec:protocol}

Every rule added after the evaluation of Section~\ref{sec:current} passed four gates before it
was kept, on shallow clones of eight repositories (ros2/demos, navigation2, moveit2,
ros2\_control\_demos, ros2\_controllers, turtlebot3, slam\_toolbox, robot\_localization):

\begin{enumerate}
\item \textbf{Precision on clean code.} The rule runs on all eight repositories; every finding is
  reviewed by hand, and a rule with a false positive is narrowed (or made silent where the fact
  cannot be known statically) until none remains.
\item \textbf{Planted fault in a real file.} The rule's fault is planted in a copy of a real
  repository file (e.g.\ a typo in the \code{node\_names} of Nav2's
  \code{map\_saver\_server.launch.py}), and the finding must appear at the planted line.
\item \textbf{Sandbox and differential test.} The fault is added to the multi-package sandbox
  workspace; every agent must trigger on it, stay off without its trigger, report only its own
  rules when selected alone, and give identical results incrementally and in full
  (Section~\ref{sec:incremental}).
\item \textbf{Message stability.} Messages contain no ROS distribution names, so fingerprints
  (baselines, \code{rosfix migrate}) do not change with \code{ROS\_DISTRO}.
\end{enumerate}

Table~\ref{tab:protocol} summarises the outcome. Gate 1 caught real problems in the checks
themselves: tf2 accepts a leading slash on published transforms but rejects it in lookups (F003
was restricted to lookup frames); rclpy runs an action server's execute callback as an executor
task (E002 is not reported under a multi-threaded executor); parameters in Python launch files
may be unevaluated dictionaries, unlike XML values (P003). It also exposed extractor gaps that
the new checks depend on: \code{PathSubstitution(...) / "x"} paths, \code{ThisLaunchFileDir()},
imports in Python launch files, plugin descriptions exported from CMake, and C++ files that use
\code{nav2\_util::LifecycleNode} without mentioning \code{rclcpp}.

\begin{table}[t]
\centering
\caption{New checks: findings on eight clean repositories and the planted fault each detects}
\label{tab:protocol}
\footnotesize
\begin{tabular}{@{}l r p{4.2cm}@{}}
\toprule
Rules & Clean & Planted fault (real file) \\
\midrule
F001--F004 & 0 & static transforms and frame ids \\
M001--M002 & 0 & Nav2 \code{map\_saver\_server} \code{node\_names} \\
C007--C008 & 0 & spawner typo, ros2\_control\_demos example 2 \\
C009 & 0 & \code{wheel\_separation: -0.10} vs.\ \code{gt<>[0.0]} \\
C010 & 1$^a$ & controller rate above the manager's \\
A001 & 0 & \code{use\_sim\_tme} in slam\_toolbox \\
P006 & 0 & descriptor range and spec validators \\
P007 & 0 & TurtleBot3 Humble Nav2 file analysed as Jazzy \\
L001, L006 & 0 & demos \code{composition} \\
L007--L008 & 0 & \code{nav2\_msgs} \\
L009 & 0 & \code{robot\_state\_publsher} in TurtleBot3 \\
L010 & 0 & component library removed from \code{install} \\
D001--D002 & 0 & domain id 233; RMW typo \\
N008 & 4$^b$ & AMCL \code{robot\_model\_type: differential} \\
N009--N010 & 1$^c$ & \code{ComputePathToPos} in Nav2's tree \\
\bottomrule
\end{tabular}

\smallskip\raggedright\footnotesize
$^a$ ros2\_control\_demos example 1: \code{joint\_state\_broadcaster} asks for 50\,Hz under a
10\,Hz controller manager. $^b$ TurtleBot3's \code{param/humble/*.yaml} analysed as Jazzy:
\code{nav2\_navfn\_planner/NavfnPlanner} was renamed \code{nav2\_navfn\_planner::NavfnPlanner}
(correct for the requested distribution). $^c$ navigation2 \code{opennav\_docking\_bt}
\code{application\_example.xml}: \code{<inverter>} instead of \code{<Inverter>}; BehaviorTree.CPP
is case-sensitive, so the tree does not load.
\end{table}

Besides the three problems noted in Table~\ref{tab:protocol}, the checks found two default
configuration files referenced by ros2\_control\_demos example 15 that do not exist (L005,
visible after the \code{PathSubstitution} fix).

\paragraph{Replaying real fixes.} \code{rosfix-history} replays fix commits: for every non-merge
commit whose message reads like a fix and which only modifies configuration files, the parent is
analysed, then the same tree with the commit's versions of those files supplied in memory; a fix
counts as detected when a finding in a modified package disappears. Matching by package rather
than by file matters, because a fix often edits \code{package.xml} while the finding is reported
where the dependency is used. On six repositories (224 candidate commits) the replay detected 9
fixes, each consistent with its commit message (six missing dependencies, an uninstalled launch
directory, a controller never spawned, a parameter typo). The candidate filter admits every
fix-worded configuration change, so this is a lower bound; recall on labelled commits is
measured by the larger study reported separately.

\subsection{Incremental Re-analysis}
\label{sec:incremental-eval}

\code{rosfix-bench --ablation} re-analyses each mutation of the benchmark three ways (full; parse
cache only; cache with agent and launch-graph reuse) and compares the findings.
Table~\ref{tab:ablation} gives median wall times. All 54 re-analyses gave identical findings. The
parse cache accounts for most of the gain on small workspaces, where resolution and agents take
0--30\,ms; launch-graph reuse matters on large ones: on navigation2 (732 files, 50 launch graphs,
1,161 node instances) the median resolution time after an edit fell from 531 to 242\,ms for a
parameter file read by most graphs and from 422 to 88\,ms for a launch file read by one graph.

\begin{table}[t]
\centering
\caption{Median re-analysis time after one mutation (ms)}
\label{tab:ablation}
\footnotesize
\begin{tabular}{@{}l r r r r@{}}
\toprule
Workspace & n & Full & Parse cache & + reuse \\
\midrule
ros2\_control\_demos & 17 & 192 & 101 & 108 \\
ros2/demos & 16 & 197 & 59 & 61 \\
turtlebot3 & 21 & 69 & 28 & 30 \\
\bottomrule
\end{tabular}
\end{table}

\subsection{Installed Packages}

With the Jazzy knowledge pack, rosfix checks a workspace that launches Nav2's installed
controller server without Nav2's source being present. \code{controller\_frequency: 20} is
reported as a type error that makes the node fail at startup (fix: \code{20.0}), and a misspelled
\code{min\_x\_velocity\_treshold} is reported as undeclared.

\textbf{Rebuilding the packs.} Rebuilding the packs from the current upstream branches
(Table~\ref{tab:packs}) exposed two defects in the shipped packs. First, endpoints that request
the system-default QoS profile were stored as if they used the ROS default profile (0 such
endpoints in the old packs, 4, 6 and 8 now), which could make the QoS agent assume a reliability
the middleware does not guarantee. Second, Nav2's AMCL declares all of its parameters through
\code{nav2\_util::LifecycleNode::add\_parameter}, a member wrapper over
\code{declare\_parameter} that the C++ extractor did not recognise. The AMCL node in the packs
had 0 (Humble), 7 (Jazzy) and 8 (Kilted) declared parameters instead of 45, 53 and 54, so every
AMCL entry of a parameter file was reported as undeclared (P003) when Nav2 was installed from
binaries, and 37 times in Nav2's own repository (e.g.\ \code{nav2\_bringup/params/nav2\_params.yaml}). The extractor now treats
\code{add\_parameter(name, default, \dots)} as a declaration; on the 12 repositories, P003
findings fell from 46 to 9. Six of the 9, in the MoveIt Servo demos, were false positives:
\code{update\_period} and \code{planning\_group\_name} are declared by moveit\_core's
acceleration filter, a plugin that \code{moveit\_servo} loads into the demo node, through a
generate\_parameter\_library listener (\code{ParamListener(node)}) whose specification lives in
a CMake subdirectory. rosfix now reads parameter-library specifications from subdirectory
\code{CMakeLists.txt} files and accepts the specifications of a node's workspace dependencies as
declarations, as it already did for their literal declarations. The remaining 3 findings, in
rmf\_ros2, are real: the composition launch file sets \code{subns}, \code{verbose} and
\code{run\_composition\_command}, which the node never declares.

We also reviewed by hand the findings of the rules added since Section~\ref{sec:current} that
fire on the 12 repositories. All 31 ROS 1 leftovers (O001, O002) are accurate: ROS 1 launch files
in the ament packages of PX4-Autopilot (6) and moveit2 (5), and 20 moveit2 C++ and Python files
that include \code{ros/ros.h} or import \code{rospy}/\code{roslib}. None of the 20 files is part
of a CMake target, so they are dead code rather than build failures. The one U002 finding is
also accurate: MoveIt's PRBT test robot gives \code{prbt\_flange}, the child of a revolute joint,
a mass of 0.02\,kg and an all-zero inertia tensor.

\subsection{Coverage}
\label{sec:coverage}

In navigation2, rosfix built 50 launch graphs with 1,412 node instances and linked 1,316 of them
to their source. Of the rest, 87 run executables from packages outside the workspace and 9 could
not be found. Parameter checks apply fully to 379 linked nodes and to top-level names only for 881
nodes that host plugins. They are disabled for 56 nodes whose declarations cannot be determined.

% ===========================================================================
\section{Case Study: TurtleBot3}

The earlier version of rosfix found three kinds of problems in the TurtleBot3 educational
platform:
\begin{itemize}
\item \textbf{Physical safety:} 6 continuous URDF joints without velocity and effort limits,
  which can damage hardware in simulation or on the real robot.
\item \textbf{Executor deadlocks:} 16 blocking \code{sleep()} calls inside ROS 2 callbacks that
  freeze single-threaded executors.
\item \textbf{Policy violations:} several launch files without \code{respawn="true"} on motor
  control nodes.
\end{itemize}
We confirmed these issues by manual inspection and testing.

The current implementation reports seven findings for TurtleBot3 (Table~\ref{tab:current}): two
sleeps inside action-server callbacks of \code{turtlebot3\_example}, and five undeclared
dependencies (\code{rclcpp\_components}, \code{camera\_ros} and \code{image\_view} in
\code{turtlebot3\_bringup}, and \code{rviz2} in the cartographer and navigation2 packages).

% ===========================================================================
\section{Related Work}

QoS Guard~\cite{lee2025qosguard} analyzes DDS QoS profiles in detail but reads only XML files, so
it cannot see QoS that source code overrides. LiSA4ROS2~\cite{zanatta2024lisa4ros2} uses abstract
interpretation on source code to synthesize SROS2 policies; it generates policies rather than
validating configuration. HAROS~\cite{santos2021haros} introduced static analysis for ROS 1 and
would need to be adapted to the DDS-based architecture and QoS model of ROS 2.
ROSDiscover~\cite{timperley2022rosdiscover} recovers and checks the run-time architecture of ROS 1
systems statically. rosfix resolves ROS 2 launch graphs and adds parameter-file binding,
installed packages and packaging checks. None of these tools compares configuration across file
types, which is the gap rosfix addresses.

The empirical studies of ROS misconfiguration~\cite{canelas2024misconfig} and interaction
bugs~\cite{chen2025ibugs} guided our choice of rules. In program repair, small models are usually
fine-tuned to generate patches~\cite{silva2023repairllama} and improve when given retrieved
examples~\cite{wang2023rapgen}. The rosfix assistant instead ranks a closed set of valid values
and verifies every choice by re-analysis.

% ===========================================================================
\section{Limitations}

\begin{description}
\item[Static analysis.] rosfix cannot detect configuration problems that appear only at runtime,
  such as parameters changed while the system runs.
\item[Dependency resolution.] The earlier version handles direct dependencies, and complex
  transitive relationships may be missed.
\item[Ground truth.] The manual validation of the earlier version covers a sample of 312 issues,
  not all 2,405.
\item[Platform.] All measurements were taken on one Windows machine.
\item[Interpreters.] The Python launch and xacro interpreters cover common subsets. Other
  constructs are marked partial, and checks that depend on them are skipped.
\item[Synthetic mutations.] The mutation operators were written together with the tool. The
  benchmark shows that the checks work on real code layouts. It does not measure recall on faults
  that occur in practice.
\item[Plugins.] Parameters that a dependency's plugin declares on its hosting node are accepted
  when the dependency is in the workspace. For a binary-installed dependency, only the knowledge
  packs' nodes are known, not the parameter libraries of their plugins.
\item[Packs and assistant.] Knowledge packs reflect the branches they were built from. The
  learned ranker was trained and evaluated on synthetic misspellings, not on fixes from real
  project histories. The sandboxed build check was not part of the evaluation because it needs
  Docker or colcon.
\end{description}

% ===========================================================================
\section{Conclusion}

We defined Shadow Configuration Drift, a class of configuration errors in ROS 2 that only appear
when files in different languages are read together, and presented rosfix to detect it. On
38,917 files from 12 repositories, rosfix reported 2,405 issues with 87.3\% precision, and
13.0\% of them could not have been found by single-file analysis. ROS 2 configuration is written
in four languages, and our results indicate that it has to be checked across all four.

The current implementation adds a resolved launch graph, knowledge packs for installed packages
and verified repair, so that rosfix can run in continuous integration on ROS 2 workspaces it has
not seen before. Future work includes dynamic analysis for runtime-only problems, learned models
for more complex drift patterns, and evaluation in more robotics domains. We also plan a
benchmark of real fixes mined from permissively licensed repositories, and to train the learned
ranker on it.

\textbf{Availability.} rosfix is open source at
\url{https://github.com/SohamWalam11/ros2-safety-lint}, together with the experimental datasets
and the benchmarks. The repair ranker is published at
\url{https://huggingface.co/SohamW11/rosfix-assist} and the knowledge packs and rosdep key lists
at \url{https://huggingface.co/datasets/SohamW11/rosfix-knowledge-packs}. The raw outputs of the
re-run benchmark and training are in \code{paper/data/}.

\bibliographystyle{IEEEtran}
\bibliography{references}

\end{document}
